\documentclass[11pt,a4paper]{article}
\usepackage[T1]{fontenc}
\usepackage[utf8]{inputenc}
\usepackage{lmodern,microtype,amsmath,amssymb,amsthm,booktabs}
\usepackage[margin=28mm]{geometry}
\usepackage[hidelinks]{hyperref}
\newtheorem{theorem}{Proposition}
\newtheorem{remark}[theorem]{Remark}
\newcommand{\E}{\mathbb E}
\newcommand{\Pp}{\mathbb P}
\newcommand{\R}{\mathbb R}
\newcommand{\pos}[1]{(#1)_+}
\setlength{\emergencystretch}{2em}
\title{Elementary reductions for the fixed-target argument}
\author{Bartosz Ko{\l}odziejek \and Kacper Rodziewicz}
\date{28 October 2026}
\begin{document}
\maketitle
\noindent These are scoped analytic reductions, not a proof of the sharp
iid Berry--Esseen conjecture. They were independently derived from the
external suggestions supplied for the present revision. AI systems were
used substantially in the derivation, checking and writing of this note.
The current manuscript and its certificates are unchanged.

\section{One fifth-moment inequality}
Let $X$ have mean zero and variance one, and assume $\E|X|^5<\infty$.
Write
\begin{gather*}
 m=\E X_+=\E X_->0,\quad \rho=\E|X|^3,\quad M_j=\E|X|^j,\\
 \beta=\E[X|X|],\quad\gamma=\E X^3,\quad\delta=\gamma-\rho\beta.
\end{gather*}
Both strict signs have positive probability, so $|\beta|<1$.
An atom at zero is allowed. Put
\[
 P_s=\frac{(\rho+\gamma)\pos{\delta}^2}{8(1+\beta)^2}
       +\frac{(\rho-\gamma)\pos{-\delta}^2}{8(1-\beta)^2}.
\]
\begin{theorem}\label{v10:primitive:fifth}
Under these hypotheses,
\begin{equation}\label{v10:primitive:combined}
 \frac{\rho M_4}{2}-\frac{M_5}{4}
 \le \frac{\rho^3}{4}-P_s-\frac{\delta^2}{16m}.
\end{equation}
Its loss dominates each of the three former losses
\[
 P_s,\qquad \frac{\delta^4}{4\rho(1-\beta^2)^2},\qquad
 \frac{\delta^2}{8m}.
\]
\end{theorem}
\begin{proof}
For one sign let $Y=X_+$ or $X_-$ and put
$a=\E Y^2$, $b=\E Y^3$, $d=b-\rho a$, and
$V=\E[Y^3(Y-\rho)^2]$. Cauchy gives
\[
 d^2=\{\E[Y^{1/2}Y^{3/2}(Y-\rho)]\}^2\le mV.
\]
If $d\ge0$, the fourth and fifth moments $c=\E Y^4$ and
$e=\E Y^5$ satisfy $c\ge b^2/a$ and $e\ge c^2/b$. Hence
\[
 V=e-2\rho c+\rho^2b\ge(c-\rho b)^2/b
   \ge b(b/a-\rho)^2=bd^2/a^2.
\]
The second inequality uses $c\ge b^2/a\ge\rho b$; it would not
follow by squaring a negative lower bound.

On the two signs,
\[
 a_\pm=(1\pm\beta)/2,\quad b_\pm=(\rho\pm\gamma)/2,
 \quad d_\pm=\pm\delta/2.
\]
Use the stronger bound on the sign with $d_\pm\ge0$, and the bound
$V_\mp\ge d_\mp^2/m$ on the other sign. Since
\[
 V_++V_-=M_5-2\rho M_4+\rho^3,
\]
division by four proves \eqref{v10:primitive:combined}. If $\delta=0$
both losses vanish and the same identity proves the claim.

For domination, suppose first that $\delta>0$. Signwise Cauchy gives
$mb_+\ge a_+^2$, so $P_s=b_+\delta^2/(16a_+^2)
\ge\delta^2/(16m)$. Thus the combined loss dominates $\delta^2/(8m)$.
Also
\[
 \delta/2=b_+a_--b_-a_+\le b_+a_-,\qquad b_+\le\rho,
\]
so $\delta^2\le4\rho b_+a_-^2$. As $1-\beta^2=4a_+a_-$,
this proves $P_s\ge\delta^4/[4\rho(1-\beta^2)^2]$.
The case $\delta<0$ is the same calculation with the signs exchanged.
\end{proof}

The new addition charges different signs of one nonnegative square.
It does not justify adding the old full weighted loss to $P_s$.
For example, the variance-one centered law
\begin{gather*}
 \Pp(X=-1/4)=396/425,\quad\Pp(X=4)=99/1700,\\
 \Pp(X=-1/2)=1/125,\quad\Pp(X=2)=1/500
\end{gather*}
violates the assertion with loss $P_s+\delta^2/(8m)$: its excess
over the available square is greater than $0.0267$. The new loss is
smaller than that square by more than $0.0201$.

\section{The first moment from one determinant}
\begin{theorem}\label{v10:primitive:first}
Suppose additionally $|X|\le R$. Necessarily $1\le\rho\le R$ and
\begin{equation}\label{v10:primitive:root}
 2m\le U(R,\rho):=
 \frac{R+\rho+\sqrt{(R-\rho)(4R^3-3R-\rho)}}{2R^2}.
\end{equation}
This upper bound equals
\begin{equation}\label{v10:primitive:inf}
 \inf_{a>0}\frac{R-\rho+2a+Ra^2}{a(2R+a)}.
\end{equation}
It is nondecreasing in $R$ and nonincreasing in $\rho$ on
$1\le\rho\le R$. Therefore a closed parameter cell may use
$U(R_{\rm upper},\rho_{\rm lower})$, with its square root rounded upwards.
\end{theorem}
\begin{proof}
For $Y=|X|$ and $u=\E Y$, the matrix
\[
 \E\left[(R-Y)\begin{pmatrix}1\\Y\end{pmatrix}
                     \begin{pmatrix}1&Y\end{pmatrix}\right]
 =\begin{pmatrix}R-u&Ru-1\\Ru-1&R-\rho\end{pmatrix}
\]
is positive semidefinite. Its determinant gives
\[
 -R^2u^2+(R+\rho)u+R^2-R\rho-1\ge0.
\]
The discriminant is $(R-\rho)(4R^3-3R-\rho)$; solving the quadratic
proves \eqref{v10:primitive:root}. The radicand is nonnegative, since
$R\ge\rho\ge1$ and $4R^3-3R-\rho\ge4R(R^2-1)\ge0$.

For completeness, the infimum identity can also be checked without a
formal minimization. For $1< R$ and $1\le\rho<R$, the upper root $U$
satisfies $1/R<U\le1<R$. Indeed the determinant polynomial is positive
at $u=1/R$, nonpositive at $u=1$, and decreases for $u\ge1$.
At
\[
 a_*=(RU-1)/(R-U)>0
\]
the quadratic
\[
 (R-U)a^2+2(1-RU)a+(R-\rho)
\]
has zero minimum, because its determinant vanishes. It is nonnegative
for every $a$. Dividing by $a(2R+a)>0$ proves that the expression in
\eqref{v10:primitive:inf} is at least $U$, with equality at $a_*$.
When $\rho=R>1$, its infimum as $a\downarrow0$ is $1/R=U$.
When $R=1$, also $\rho=1$ and the expression is identically one.

For fixed $a>0$ the expression decreases in $\rho$. Its derivative
in $R$ has the sign of
$a^3-3a+2\rho\ge-2+2\rho\ge0$.
Taking infima proves both monotonicity statements, including endpoints.
\end{proof}

This replaces the eleven sampled auxiliary values of $a$ by their exact
optimal value. It does not replace the stronger signwise matrix constraints;
those remain separate valid information about the same law.

\section{A calculus proof for the cosine constant}
\begin{theorem}\label{v10:primitive:cosine}
For every $0\le x\le1$,
\[
 1-\cos(2\pi x)\le\frac{4553}{1000}x.
\]
\end{theorem}
\begin{proof}
For $y>0$ put $h(y)=(1-\cos y)/y$. The derivative has the sign of
$N(y)=y\sin y+\cos y-1$. On $(0,\pi)$ its sign is the sign of
$y-\tan(y/2)$. The function $\tan(y/2)-y$ decreases up to $\pi/2$
and increases thereafter to infinity; it is negative at $\pi/2$.
It therefore has exactly one zero $y_*$ in $(\pi/2,\pi)$.
On $(\pi,2\pi)$, $N(y)<0$. Thus the maximum of $h$ on $[0,2\pi]$
is attained at $y_*$, and the half-angle identity gives
\[
 h(y_*)=\frac{2y_*}{1+y_*^2}.
\]

Take $y_0=23311/10000$. The six-term alternating lower bound is
\begin{equation}\label{v10:primitive:cosine-check}
 N(y_0)>\sum_{j=1}^{6}
   \frac{(-1)^{j-1}(2j-1)y_0^{2j}}{(2j)!}
 >\frac3{200000}>0.
\end{equation}
The alternating remainder has the asserted sign: the ratio of successive
term magnitudes is $y_0^2/[(2j-1)(2j+2)]$, which is below one for every
$j\ge6$ and decreases. Since $\pi/2<y_0<\pi$, this proves $y_* > y_0$.
The function $y/(1+y^2)$ decreases for $y>1$. Using $\pi<355/113$, we get
\[
 \max_{0<x\le1}\frac{1-\cos(2\pi x)}x
 =\frac{4\pi y_*}{1+y_*^2}
 <\frac{4(355/113)y_0}{1+y_0^2}
 <\frac{4553}{1000}-\frac1{10000}.
\]
The last comparison and \eqref{v10:primitive:cosine-check} are exact
rational inequalities. To make the elementary constant enclosure
reproducible, use
$\pi=16\arctan(1/5)-4\arctan(1/239)$ and the alternating arctangent
series, with eight and two terms respectively. Their outward intervals
already prove $3<\pi<355/113$. Machin's identity follows by the tangent
addition formula and the fact that the angle in question lies in $(0,\pi/2)$.
The case $x=0$ is immediate.
\end{proof}

The constant $23/5$ follows since $4553/1000<23/5$.
Consequently both former Bernstein coefficient tables can be replaced by
this one calculus argument and its few rational comparisons. This changes
neither the harmonic span domain nor the other count-comparison hypotheses.
The value $9/2$ is not a valid replacement constant: at $x=3/8$ the ratio
is $(8+4\sqrt2)/3>9/2$.

\section{A common two-sign transfer for every sample size}
Here only $\E X=0$, $\E X^2=1$ and $\rho=\E|X|^3<\infty$ are needed;
bounded support is not required. Define probability measures on $(0,\infty)$ by
\[
 \mu_\pm(dy)=\frac y m\Pp(\pm X\in dy,\ \pm X>0),
 \qquad c_\pm=\int y\,d\mu_\pm(y)>0.
\]
These are the distributions obtained by weighting each strict sign by its
absolute magnitude. Define
\[
 Q=\frac m{c_-}\delta_{-c_-}
  +\left(1-\frac m{c_-}-\frac m{c_+}\right)\delta_0
  +\frac m{c_+}\delta_{c_+},\qquad
 \eta=\rho(P)-\rho(Q).
\]
\begin{theorem}\label{v10:primitive:transfer}
$Q$ is a probability measure with the same first two moments as $P$,
and
\begin{equation}\label{v10:primitive:variance}
 \eta=m\{\operatorname{Var}_{\mu_-}Y+
                  \operatorname{Var}_{\mu_+}Y\}\ge0.
\end{equation}
For the characteristic functions $\psi$ of $P$ and $g$ of $Q$,
\begin{equation}\label{v10:primitive:char}
 |\psi(t)-g(t)|\le\frac{|t|^3}{6}\eta\qquad(t\in\R).
\end{equation}
For every integer $n\ge1$,
\begin{equation}\label{v10:primitive:power}
 |\psi(t)^n-g(t)^n-ng(t)^{n-1}(\psi(t)-g(t))|
 \le\frac{n(n-1)}2|\psi(t)-g(t)|^2
 \le\frac{n(n-1)}{72}|t|^6\eta^2.
\end{equation}
All statements include zero atoms and $\eta=0$.
\end{theorem}
\begin{proof}
On each sign Cauchy gives
$m^2\le\Pp(\pm X>0)\E X_\pm^2
=\Pp(\pm X>0)mc_\pm$.
Thus $m/c_\pm\le\Pp(\pm X>0)$ and the mass at zero in $Q$ is
nonnegative. Its strict-sign first and second moments equal those of $P$:
they are $m$ and $mc_\pm$. Its third absolute moment is
$m(c_-^2+c_+^2)$, whereas that of $P$ is
$m\int y^2(d\mu_-+d\mu_+)$. This proves \eqref{v10:primitive:variance}.
If $\eta=0$, both $\mu_\pm$ are point masses at $c_\pm$; the original
law is then exactly $Q$, including its possible zero atom.

For $\sigma\in\{-1,1\}$ put
\[
 h_{\sigma,t}(y)=\frac{e^{i\sigma ty}-1}{y}
      =i\sigma t\int_0^1 e^{i\sigma tsy}\,ds,
\]
with the continuous definition at zero. The integral representation gives
\[
 h_{\sigma,t}''(y)=(i\sigma t)^3\int_0^1s^2e^{i\sigma tsy}\,ds,
 \qquad |h_{\sigma,t}''(y)|\le |t|^3/3.
\]
Taylor's formula along the real segment from $c_\sigma$ to $y$ yields
\[
 h(y)-h(c_\sigma)-h'(c_\sigma)(y-c_\sigma)
 =(y-c_\sigma)^2\int_0^1(1-r)
             h''(c_\sigma+r(y-c_\sigma))\,dr.
\]
Every integral is justified by the finite second moment of $\mu_\sigma$.
The linear term integrates to zero. Finally,
\[
 \psi(t)-g(t)=m\sum_{\sigma\in\{-1,1\}}
 \left\{\int h_{\sigma,t}(y)\,d\mu_\sigma(y)
                         -h_{\sigma,t}(c_\sigma)\right\}.
\]
Taking absolute values proves \eqref{v10:primitive:char}. The same two
measures occur at every frequency; no independent frequencywise law has
been introduced.

For $n\ge2$, the segment joining the two characteristic functions lies
in the closed unit disk. Taylor's formula for $z^n$ along this segment
bounds the second remainder by $n(n-1)|\psi-g|^2/2$.
For $n=1$ the remainder is zero. This proves \eqref{v10:primitive:power}.
\end{proof}

For the normalized sum, replace $t$ by $t/\sqrt n$ in
\eqref{v10:primitive:power}; the last bound becomes
$(n-1)|t|^6\eta^2/(72n^2)$. This is an all-$n$ transfer, not a
uniform CDF estimate: inversion, smoothing reserves and any signed
first-order improvement remain separate requirements.

The coefficient $1/6$ in \eqref{v10:primitive:char} is sharp. For
$0<\varepsilon<1/6$, take masses
$1/2-\varepsilon,3\varepsilon,1/2-3\varepsilon,\varepsilon$
at $-1,0,1,2$. This law is centered with variance one. Only the positive
size-biased sign has nonzero variance, so
$\E_P X^3-\E_Q X^3=\eta>0$. Its characteristic difference has leading
term $-it^3\eta/6$, with absolute remainder at most
$|t|^4(\E_P X^4+\E_Q X^4)/24$. Thus the proposed smaller coefficient
$1/7$ fails whenever
$0<|t|<4\eta/[7(\E_P X^4+\E_Q X^4)]$.

\begin{remark}[What the transfer does not replace]
The reference $Q$ can have three atoms and is not automatically Bernoulli.
The modulus bound alone cannot supply a favorable signed Fourier response,
a sharp finite-base comparison, or the stronger zero-mass radius used in
some existing numerical estimates. The earlier response and price
identities must remain wherever those extra constraints are used.
The present formulation replaces only the scalar defect/response/remainder
construction and supplies a common version for every $n$.
\end{remark}

\section{One compact domain for the finite base}
This last reduction uses an external inequality, unlike the elementary
moment and transfer propositions above. Set
$c=(3+\sqrt{10})/(6\sqrt{2\pi})$ and let $\Delta_n(P)$ denote the
Kolmogorov distance from the normalized iid sum to the standard normal.
\begin{theorem}\label{v10:primitive:finite}
Assume the universal affine bound
\[
 \sqrt n\,\Delta_n(P)\le\frac{3031}{10000}
                         \left(\rho+\frac{323}{500}\right).
\]
For $1\le n\le41$, every attaining fixed-target pair $(P,H)$ with
\[
 a=\sqrt n\{\Pp(S_n\le H)-\Phi(H/\sqrt n)\}-c\rho\ge0
\]
satisfies
\[
 \rho<\frac{7361}{4000}<1.841,\qquad
 \operatorname{supp}P\subset(-13,13),\qquad |H|<26.
\]
No predecessor inequality is assumed.
\end{theorem}
\begin{proof}
The moment bound follows from $c>819/2000$ and direct division:
\[
 \rho<\frac{(3031/10000)(323/500)}{819/2000-3031/10000}
      =\frac{7361}{4000}.
\]
The fixed-target contamination inequality is
\begin{align*}
 I(y)={}&n^{3/2}(h(y)-F)+ny\phi(H/\sqrt n)
       +\tfrac12H\phi(H/\sqrt n)(y^2-1)\\
       &+\tfrac32c\rho(y^2-1)-c(|y|^3-\rho-3y\beta)\le0,
\end{align*}
with equality on the closed support, where
$F=\Pp(S_n\le H)$ and $h(y)=\Pp(S_{n-1}\le H-y)$.
For $n=1$ the same contamination calculation has
$h(y)=\mathbf1_{\{y\le H\}}$; no deleted random summand is required.
The usual fixed-$n$ attainment argument also applies at $n=1$: its tightness,
moving closed threshold and third-moment penalty steps do not require
$n\ge2$. Thus the finite-base extremal reduction has no exceptional first
index.

Let $u=|y|\ge1$. We have $|\beta|\le1$,
$\phi(0)<2/5$, $\sup_z|z\phi(z)|=\phi(1)<1/4$,
$\sqrt n<13/2$, and $819/2000<c<41/100$.
In particular $n^{3/2}<533/2$ and
\[
 I(y)\le P(u):=266.5+1.944715u^2+17.63u-0.4095u^3.
\]
To check the coefficients, the linear coefficient is
$41(2/5)+3(41/100)=17.63$ and the quadratic coefficient is
$13/16+(3/2)(41/100)(1841/1000)=1.944715$.
The constant contribution of the two $c\rho$ terms is $-c\rho/2$,
which can be dropped; there is no additional positive constant to retain.
The term with $H\phi$ is bounded by $(13/16)(u^2-1)$ before its
negative constant is dropped. All displayed decimals are exact rationals.

Exact evaluation gives
\[
 P(13)<-75,\quad P'(13)<-139,\quad P''(13)<-28,
 \quad P'''=-2457/1000<0.
\]
It follows successively that $P'',P',P$ are negative on $[13,\infty)$.
Closed-support contact therefore excludes both exterior rays, including
their boundary points.

If $H<0$, Chebyshev gives $F\le n/H^2$; if $H>0$, the positive
discrepancy is at most $1-\Phi(H/\sqrt n)\le n/H^2$.
Thus $|H|\ge26$ would imply
\[
 \sqrt n(F-\Phi(H/\sqrt n))<\frac{533/2}{26^2}
   =\frac{819}{2000}-\frac{397}{26000}<c\rho,
\]
contradicting $a\ge0$.
The elementary constants used above follow from the same directed
$\pi$ bounds as in the cosine proof, directed squares around $\sqrt{10}$,
and $e>8/3$, which gives $\phi(1)<1/4$ from $\pi>3$.
\end{proof}

This is one compact finite-index problem, not a solution of that problem.
The reference law in Proposition~\ref{v10:primitive:transfer} can still
have three atoms, and a complete estimate of its sums and the original-law
response remains necessary.

\section{Verification boundary}
The proofs above establish their stated analytic implications. The accompanying
exact script checks rational coefficient identities, scalar signs and a
fixed adversarial collection of actual centered variance-one laws, including
atoms at zero, reflected skew laws, arbitrarily small mixture weights and
a distant atom. It also rejects the false stronger fifth-moment sum and
the false cosine and characteristic constants described above. Such tests
are checks on the derivations, not substitutes for their proofs.

These reductions need no new dense parameter cover. Applying their stronger
bounds to the full entry or local comparison still requires the original
domain and coverage conditions. No old proof object is deleted, no target
is weakened, and no complete finite-base or global conclusion is asserted.
\end{document}
